\chapter{Templates, Checklists, and Ready-to-Use Resources}

\section{Chapter Overview}
This chapter gathers practical templates, checklists, and prompts from earlier sections so teams can adopt collaboration practices without starting from scratch. Customize each template for your context; treat them as living documents that evolve with your process.

\section{Commit Message Template}
\begin{lstlisting}[style=shell]
type(scope): brief summary (50 chars max)

Context:
- why was this needed?
- what alternatives were considered?

Testing:
- unit:
- integration:
- manual:

Follow-up:
- feature flag plan or cleanup steps

Refs: ticket/incident IDs
\end{lstlisting}

\section{PR Description Template}
\begin{lstlisting}[style=markdown]
## Summary
[Brief description of changes]

## Motivation
[Why this change is needed and who it impacts]

## Implementation Details
[Key technical decisions, trade-offs, and diagrams/links]

## Testing
- [ ] Unit tests
- [ ] Integration tests
- [ ] Manual QA (steps recorded below)

### Evidence
- screenshots/videos/logs
- monitoring links

## Rollout Plan
- default flag values, canary plan, rollback steps

## Checklist
- [ ] Docs updated or N/A
- [ ] Migrations reversible and tested
- [ ] Breaking changes communicated
- [ ] Observability (metrics/logs) in place
\end{lstlisting}

\section{Review Checklist Card}
\begin{enumerate}
  \item \textbf{Intent}: Does the PR description explain motivation and scope?
  \item \textbf{Correctness}: Do tests and logic cover happy and failure paths?
  \item \textbf{Security/Privacy}: Any injection, data leak, or compliance risk?
  \item \textbf{Performance}: Are there obvious hot-path regressions?
  \item \textbf{Maintainability}: Is the code readable, documented, and consistent with architecture?
  \item \textbf{Testing Evidence}: Do screenshots/logs prove the change works?
  \item \textbf{Rollout Safety}: Are feature flags, migrations, or rollback plans documented?
\end{enumerate}

\section{Collaboration Retro Agenda}
\begin{enumerate}
  \item \textbf{Warm-up}: Celebrate one collaboration win from the past sprint.
  \item \textbf{Metrics Review}: Cycle time, review latency, automation failures.
  \item \textbf{Start/Stop/Continue}: Collect ideas on commits, PRs, reviews.
  \item \textbf{Experiment Selection}: Choose one experiment, assign owner/due date.
  \item \textbf{Template/Bot Feedback}: Gather pain points for automation UX.
\end{enumerate}

\section{Feature Flag Lifecycle Checklist}
\begin{enumerate}
  \item Document flag purpose, owner, and default value.
  \item Log rollout stages (internal, canary, general) with dates.
  \item Track metrics/alerts tied to flag activation.
  \item Schedule cleanup date; add ticket for removing stale flags.
  \item After removal, confirm docs/tests are updated.
\end{enumerate}

\section*{Exercises}

\paragraph{1. Template adoption} Pick one template above and pilot it for two weeks. Survey contributors about clarity and adjust accordingly.

\paragraph{2. Checklist walk-through} Run a code review using the review checklist card. Note which items surfaced new insights versus duplicating intuition.

\paragraph{3. Retro rehearsal} Facilitate the collaboration retro agenda verbatim. Capture the experiment selected and track its outcome.

\paragraph{4. Flag lifecycle audit} Use the checklist to audit three existing feature flags. Log missing owners or cleanup dates.

\paragraph{5. Automation UX feedback} Add a ``template/bot feedback'' section to your next retro and collect at least three improvement ideas.
